\indent ವಿ. ಸೂ. :ಷಟ್ಪದಿಯ ಮೂರು ಮತ್ತು ಆರನೆಯ ಸಾಲುಗಳಲ್ಲಿನ ಕೊನೆಯ ಅಕ್ಷರ “ಲಘು” ಆಗಿದ್ದರೂ, ಅದನ್ನು ಗುರುವೆಂತಲೇ ತಿಳಿಯ ಬೇಕು.\par

\vspace{4pt plus 4pt minus 1pt}\noindent{\kbvsecfont\fontsize{11pt}{14.5pt}\selectfont\color{kbvaccentdark} (4)\hspace{0.45em}ರಗಳೆ :}\hspace{0.5em}ಇದು ಮಾತ್ರಾಗಣಕ್ಕೆ ಸೇರಿದ ಪದ್ಯ ಜಾತಿ, “ರಘಟಾ” ಎಂಬ ಸಂಸ್ಕೃತ ಶಬ್ದದ ತದ್ಭವ ರೂಪವೇ “ರಗಳೆ”. ರಗಳೆಗಳಲ್ಲಿ ಉತ್ಸಾಹ ರಗಳೆ, ಮಂದಾನಿಲ ರಗಳೆ, ಲಲಿತ ರಗಳೆ ಎಂಬ ಮೂರು ಪ್ರಕಾರ ಗಳು ಮುಖ್ಯವಾದವುಗಳು. ರಗಳೆಗಳಲ್ಲಿ ಇಂತಿಷ್ಟು ಪಾದಗಳಿರಬೇಕೆಂಬ ನಿಯಮವಿಲ್ಲ; ಆದರೆ ಮಾತ್ರಾಗಣಗಳು ಪ್ರತಿಯೊಂದು ಪಾದದಲ್ಲೂ ಒಂದೇ ಸಮನಾಗಿರುವುದು; ಮತ್ತು ಪ್ರತಿ ಎರಡು ಪಾದಗಳಲ್ಲಿ ಪ್ರಾಸನಿಯಮವಿರುವುದು. ಕೆಲವೊಂದು ಸಂದರ್ಭಗಳಲ್ಲಿ ಪಾದದ ಆದ್ಯಂತಗಳಲ್ಲಿ ಪ್ರಾಸವಿರುವುದೂ ಉಂಟು.\par

\indent ರಗಳೆಗಳ ಲಕ್ಷಣಗಳು\par

\vspace{4pt plus 4pt minus 1pt}\noindent{\kbvsecfont\fontsize{11pt}{14.5pt}\selectfont\color{kbvaccentdark} (1)\hspace{0.45em}ಉತ್ಸಾಹ ರಗಳೆ:}\hspace{0.5em}ಈ ಜಾತಿಯ ಪದ್ಯಗಳಲ್ಲಿ ಪ್ರತಿಯೊಂದು ( ೧೧ ೬) ಪಾದದಲ್ಲೂ ಮೂರು ಮಾತ್ರೆಯ ನಾಲ್ಕು ಮಾತ್ರಾಗಣಗಳಿರುವುವು.\par

\vspace{4pt plus 4pt minus 1pt}\noindent{\kbvsecfont\fontsize{11pt}{14.5pt}\selectfont\color{kbvaccentdark} (2)\hspace{0.45em}ಮಂದಾನಿಲ ರಗಳೆ:}\hspace{0.5em}ಈ ಜಾತಿಯ ಪದ್ಯಗಳಲ್ಲಿ, ಪ್ರತಿ ಯೊಂದು ಪಾದದಲ್ಲೂ ನಾಲ್ಕು ಮಾತ್ರೆಯ ನಾಲ್ಕು ಮಾತ್ರಾಗಣಗಳಿರು ವುವು; ಅಥವಾ ಪ್ರತಿಯೊಂದು ಪಾದದಲ್ಲೂ ಮೂರು ಮತ್ತು ಐದು\par

\indent ಮಾತ್ರೆಯ ಮಾತ್ರಾಗಣಗಳು (3 + 5 + 3 + 5) ತೆರನಾಗಿ ಸಂಭವಿಸುವುವು.\par

\vspace{4pt plus 4pt minus 1pt}\noindent{\kbvsecfont\fontsize{11pt}{14.5pt}\selectfont\color{kbvaccentdark} (3)\hspace{0.45em}ಲಲಿತ ರಗಳೆ:}\hspace{0.5em}ಈ ಜಾತಿಯ ಪದ್ಯಗಳಲ್ಲಿ ಪ್ರತಿಯೊಂದು ಪಾದದಲ್ಲೂ ಐದು ಮಾತ್ರೆಯ ನಾಲ್ಕು ಮಾತ್ರಾಗಣಗಳಿರುವುವು.\par

\indent ಸೂಚನೆ: ರಗಳೆಗಳಲ್ಲಿ ಪ್ರತಿ ಎರಡು ಪಾದಗಳಲ್ಲಿ ಆದಿ ಪ್ರಾಸವಿರು ವುದು; ಕೆಲವೊಂದು ಸಂದರ್ಭಗಳಲ್ಲಿ ಆದಿಪ್ರಾಸದ ಜೊತೆಗೆ ಅಂತ್ಯ ಪ್ರಾಸವೂ ಇರುವುದುಂಟು. ಅಪರೂಪವಾಗಿ ಅಂತ್ಯಪ್ರಾಸ ಮಾತ್ರವಿರುವ ಸಂದರ್ಭಗಳೂ ಉಂಟು.\par
